% BHU PYQ -- Manually transcribed & error-corrected
% Paper: GLB-601: Palaeontology
% Exam: B.Sc. (Hons.) Semester VI Examination, 2022-23

\documentclass[11pt,a4paper]{article}
\usepackage[a4paper,top=2.5cm,bottom=2.5cm,left=2.8cm,right=2.8cm,headheight=14pt]{geometry}
\usepackage{amsmath,amssymb}
\usepackage[T1]{fontenc}
\usepackage[utf8]{inputenc}
\usepackage{lmodern,microtype,fancyhdr,enumitem,hyperref,lastpage,parskip}

\hypersetup{colorlinks=true,urlcolor=black,linkcolor=black,
  pdftitle={GLB-601 Palaeontology -- BHU B.Sc. Sem VI 2022-23}}
\pagestyle{fancy}\fancyhf{}
\renewcommand{\headrulewidth}{0.4pt}\renewcommand{\footrulewidth}{0.4pt}
\fancyhead[L]{\small\textit{Banaras Hindu University}}
\fancyhead[R]{\small\textit{GLB-601: Palaeontology}}
\fancyfoot[C]{\small Page \thepage\ of \pageref{LastPage}}

\newlist{parts}{enumerate}{2}
\setlist[parts,1]{label=(\alph*),leftmargin=*,itemsep=5pt,topsep=3pt}
\newcommand{\pts}[1]{\hfill\textit{[#1]}}

\begin{document}
\begin{center}
  {\large\bfseries BANARAS HINDU UNIVERSITY}\\[2pt]
  {\normalsize B.Sc. (Hons.) --- Semester VI Examination, 2022-23}\\[6pt]
  \rule{0.8\linewidth}{0.5pt}\\[6pt]
  {\large\bfseries Paper: GLB-601 --- Palaeontology}\\[4pt]
  {\normalsize Time: 3 Hours \hfill Maximum Marks: 70}\\[2pt]
  {\small Ref. No.: 058440}
\end{center}
\rule{\linewidth}{0.4pt}
\smallskip
\noindent\textit{Instructions: Answer \textbf{five} questions in all, selecting at least \textbf{two} each from Sections A and B, and Section-C which is compulsory. All questions are of equal value.}
\bigskip

\bigskip\noindent{\large\bfseries SECTION --- A}
\rule{\linewidth}{0.4pt}\smallskip

\noindent\textbf{Question 1.} Discuss different techniques of collection and preparation of microfossils. \pts{14}

\medskip\noindent\textbf{Question 2.} Write short notes on \textbf{any two} of the following: \pts{$2 \times 7 = 14$}
\begin{parts}
  \item Darwinism
  \item Shell coiling in gastropods
  \item Binomial nomenclature
\end{parts}

\medskip\noindent\textbf{Question 3.} Describe the diagnostic morphological features of brachiopods. \pts{14}

\medskip\noindent\textbf{Question 4.} Write short notes on \textbf{any four} of the following: \pts{$4 \times 3\frac{1}{2} = 14$}
\begin{parts}
  \item Stratigraphic significance of conodonts
  \item Organic-walled microfossils
  \item Concept of species
  \item Sutures in ammonoids
  \item Nature of umbo in bivalves
\end{parts}

\bigskip\noindent{\large\bfseries SECTION --- B}
\rule{\linewidth}{0.4pt}\smallskip

\noindent\textbf{Question 5.} Discuss the paleoclimatic significance of plant fossils. \pts{14}

\medskip\noindent\textbf{Question 6.} Write short notes on \textbf{any two} of the following: \pts{$2 \times 7 = 14$}
\begin{parts}
  \item Ornamental features of ammonoid shells
  \item Siwalik vertebrate fauna
  \item Determination of left- and right-valves in bivalves
\end{parts}

\medskip\noindent\textbf{Question 7.} Discuss briefly \textbf{any two} of the following: \pts{$2 \times 7 = 14$}
\begin{parts}
  \item Features distinguishing brachiopods from bivalves
  \item Difference between ammonoids and nautiloids
  \item Shell forms in gastropods
\end{parts}

\bigskip\noindent{\large\bfseries SECTION --- C}
\rule{\linewidth}{0.4pt}\smallskip

\noindent\textbf{Question 8.} Define the following in a maximum of 5 sentences each: \pts{$7 \times 2 = 14$}
\begin{parts}
  \item Sinistral coiling in gastropods
  \item Lunule in bivalves
  \item Mixoperipheral growth in brachiopods
  \item Columella
  \item Delthyrium
  \item Umbilicus in ammonoids
  \item Protegulum and prodissoconch
\end{parts}

\vfill\rule{\linewidth}{0.4pt}
\begin{center}\small
\href{https://bhu-exam.vercel.app/}{Click here to visit our website}\\[4pt]
\href{https://me-aryan.vercel.app/}{Click here to visit our contributor}\\[4pt]
\href{https://quashan-me.vercel.app/}{Click here to visit our contributor}
\end{center}
\end{document}
